\apendice{Documentación técnica de programación}

\section{Introducción}
El presente apéndice conforma el Manual del Programador y de Pruebas del sistema HIVES. Su objetivo principal es documentar la arquitectura a nivel de código, la organización del repositorio, el entorno necesario y los procedimientos de validación del software seguidos. Esta documentación está orientada a facilitar la mantenibilidad, futura escalabilidad y el despliegue del sistema por parte de otros ingenieros o desarrolladores que deseen auditar o continuar el proyecto.

\section{Estructura de directorios}
El repositorio del proyecto ha sido organizado siguiendo el patrón estándar para aplicaciones modulares en Python, separando claramente la lógica de negocio, los datos, los modelos de Inteligencia Artificial y los recursos estáticos. La jerarquía principal de directorios es la siguiente:

\dirtree{%
.1 Trabajo-Fin-de-Grado-2025-2026/.
.2 arduino.
.3 LecturasParaInterfaz{.}ino\DTcomment{Firmware del sensor (C++)}.
.2 release\DTcomment{Contiene los ejecutables de Windows y Linux.}.
.2 data\DTcomment{Almacena datos y logs}.
.3 data{.}db\DTcomment{Base de datos SQLite (persistencia local)}.
.3 hives{.}log\DTcomment{Logs para diagnóstico de errores}.
.2 models\DTcomment{Lógica de entrenamiento y modelos}.
.3 mejor\_modelo{.}keras\DTcomment{Modelo en producción (TensorFlow)}.
.3 clases{.}json\DTcomment{Nombres de clase en inferencia}.
.3 mejor\_modelo\_final{.}keras\DTcomment{Modelo anterior (descartado)}.
.3 pruebaColab{.}ipynb\DTcomment{Notebook activo de entrenamiento}.
.3 pruebaTensorDataSetCorregido{.}ipynb\DTcomment{Notebook anterior (descartado)}.
.2 src/.
.3 main{.}py\DTcomment{ENTRY POINT de la aplicación}.
.3 constants{.}py.
.3 sensor{.}py.
.3 runtime\_hook{.}py.
.3 db/\DTcomment{Inicializador de la base de datos}.
.4 \_\_init\_\_{.}py.
.4 seeder{.}py.
.3 hooks/\DTcomment{Empaquetado de ML con PyInstaller}.
.4 hook-keras{.}py.
.4 hook-matplotlib{.}py.
.4 hook-tensorflow{.}py.
.3 hives/\DTcomment{CÓDIGO ACTIVO (paquete hives.*)}.
.4 \_\_init\_\_{.}py.
.4 constants{.}py\DTcomment{Constantes globales del proyecto}.
.4 core/.
.5 \_\_init\_\_{.}py.
.5 sensor{.}py\DTcomment{Comunicación serie (SerialReader)}.
.5 database{.}py\DTcomment{Inicialización del esquema de la BD}.
.5 paths{.}py\DTcomment{Variables de rutas}.
.4 gui/\DTcomment{Capa de presentación (Vista MVC)}.
.5 \_\_init\_\_{.}py.
.5 main\_window{.}py\DTcomment{Ventana principal (SpectroControlUI)}.
.5 widgets{.}py\DTcomment{Canvas espectral (SpectralCanvas)}.
.5 workers{.}py\DTcomment{Hilos QThread asíncronos}.
.5 dialogs{.}py\DTcomment{Diálogos modales}.
.4 inference/.
.5 \_\_init\_\_{.}py.
.5 model{.}py\DTcomment{Carga y ejecución del modelo}.
.4 reports/.
.5 \_\_init\_\_{.}py.
.5 pdf\_report{.}py\DTcomment{Generación de informes PDF}.
.2 tests/\DTcomment{Pruebas automatizadas (pytest)}.
.3 conftest{.}py\DTcomment{Fixtures compartidos (BD, mocks, ventana)}.
.3 test\_database{.}py.
.3 test\_dialogs{.}py.
.3 test\_inference{.}py.
.3 test\_main\_window{.}py.
.3 test\_pdf\_report{.}py.
.3 test\_sensor{.}py.
.3 test\_widgets{.}py.
.3 test\_workers{.}py.
.2 build\_linux{.}sh\DTcomment{Construcción en Linux}.
.2 build\_macos{.}sh\DTcomment{Construcción en macOS}.
.2 build\_windows{.}ps1\DTcomment{Construcción en Windows}.
.2 coverage{.}xml\DTcomment{Reporte de cobertura (pytest)}.
.2 pyproject{.}toml\DTcomment{Configuración de pytest}.
.2 requirements{.}txt\DTcomment{Dependencias Python}.
.2 sonar\_project{.}properties\DTcomment{Configuración SonarQube}.
.2 HIVES{.}spec\DTcomment{Configuración para PyInstaller}.
.2 README{.}md\DTcomment{Documentación del repositorio}.
}

\section{Manual del programador}

Este apartado detalla las dependencias y las decisiones de implementación
adoptadas durante el desarrollo del sistema.

\subsection*{Entorno de desarrollo y dependencias}

El núcleo de la aplicación ha sido desarrollado íntegramente en
\textbf{Python 3.10}, utilizando \textbf{Visual Studio Code} como Entorno de
Desarrollo Integrado (IDE) y \textbf{Git} para el control de versiones,
alojando el repositorio en GitHub.

El proyecto hace uso de las siguientes librerías de terceros:

\begin{itemize}
    \item \textbf{PyQt6:} \textit{framework} para la construcción de la
    interfaz gráfica de usuario (GUI) y la gestión de hilos mediante
    \texttt{QThread}.
    \item \textbf{TensorFlow:} biblioteca para la carga del modelo de red
    neuronal y la ejecución de la inferencia en producción.
    \item \textbf{Scikit-learn:} biblioteca empleada durante la fase de
    experimentación para el entrenamiento y evaluación de clasificadores
    alternativos, incluyendo \textbf{Random Forest}.
    \item \textbf{cuML:} biblioteca utilizada durante la fase de
    experimentación para el entrenamiento y comparativa de clasificadores
    de alto rendimiento, incluyendo nuevamente \textbf{Random Forest}.
    \item \textbf{XGBoost}: biblioteca empleada durante la fase de experimentación
    para el entrenamiento y evaluación de clasificadores alternativos, incluyendo \textbf{XGBClassifier},
    que terminó integrándose como la primera etapa de procesamiento dentro de la arquitectura híbrida definitiva.
    \item \textbf{PySerial:} módulo empleado para establecer la comunicación
    de bajo nivel con el Arduino Uno a través del puerto serie.
    \item \textbf{FPDF2:} librería para la maquetación y generación de los
    reportes en formato PDF.
\end{itemize}

Adicionalmente, se hace un uso extensivo de la librería estándar de Python
(\texttt{sqlite3}, \texttt{os}, \texttt{sys}, \texttt{time}, \texttt{datetime}).
\subsection*{Detalles de implementación: concurrencia y base de datos}

Para evitar la congelación de la interfaz gráfica (\textit{GUI freezing})
durante la lectura del puerto serie y la inferencia del modelo, se ha
implementado un sistema de concurrencia basado en \textbf{QThread} de PyQt6.
Este hilo secundario se levanta exclusivamente al iniciar el análisis y emite
una señal a la vista principal una vez obtenidos los resultados, manteniendo
la interfaz reactiva en todo momento.

En cuanto a la persistencia de datos, se ha prescindido de mapeadores
objeto-relacional (ORM) en favor de consultas SQL puras sobre SQLite. Esta
decisión minimiza el consumo de recursos y otorga un control optimizado sobre
las sentencias de inserción y filtrado del histórico. El esquema de la base de
datos, compuesto por las tablas \texttt{muestras}, \texttt{predicciones},
\texttt{analisis} y \texttt{calibraciones}, ha sido inicializado mediante un
script de \textit{inicialización} y  \textit{seeding} que garantiza la creación de la estructura relacional
en el primer arranque de la aplicación.
\newpage
\section{Pruebas del sistema}

Para asegurar la calidad, fiabilidad y robustez de la aplicación, se ha diseñado
una batería de pruebas centrada en los componentes críticos del sistema.

\subsection*{Batería de pruebas automatizadas}

El proyecto cuenta con una suite de pruebas unitarias y de integración
implementada mediante el \textit{framework} \textbf{pytest\cite{pytest2024docs}}, ejecutable
de forma completa con el siguiente comando desde la raíz del repositorio:

\begin{verbatim}
pytest tests/
\end{verbatim}

La suite se organiza en ocho módulos de test independientes, cada uno
orientado a un componente específico de la arquitectura:

\begin{itemize}
    \item \texttt{test\_database.py}: verifica la correcta creación del
    esquema de la base de datos y la integridad referencial entre las
    tablas \texttt{muestras}, \texttt{predicciones}, \texttt{analisis}
    y \texttt{calibraciones}.
    \item \texttt{test\_sensor.py}: prueba el módulo de comunicación
    serie, comprobando el comportamiento ante desconexiones abruptas,
    \textit{timeouts} y tramas de datos incompletas o corruptas.
    \item \texttt{test\_inference.py}: valida que el módulo de
    clasificación recibe tensores con la dimensionalidad y normalización
    correctas y emite predicciones coherentes con la arquitectura del
    modelo.
    \item \texttt{test\_main\_window.py}: prueba la lógica de la ventana
    principal, incluyendo el ciclo de análisis completo y la gestión
    de estados de la interfaz.
    \item \texttt{test\_widgets.py}: comprueba el funcionamiento del
    canvas espectral (\texttt{SpectralCanvas}) y su renderizado.
    \item \texttt{test\_workers.py}: verifica el comportamiento de los
    hilos \textbf{QThread}, asegurando que las señales se emiten
    correctamente al hilo principal una vez completada la operación.
    \item \texttt{test\_dialogs.py}: valida los diálogos modales de la
    aplicación y su interacción con el flujo principal.
    \item \texttt{test\_pdf\_report.py}: comprueba la generación de
    informes PDF, verificando que el documento resultante contiene los
    campos requeridos por el requisito RF-07.
\end{itemize}

\subsection*{Estrategia de \textit{mocking} y \textit{fixtures}}

Dado que el sistema interactúa con hardware externo y con dependencias
de gran peso computacional, la suite de pruebas hace un uso extensivo
de \textit{mocking\cite{fowler2007mocks}} para aislar cada componente y garantizar que los
tests son reproducibles en cualquier entorno, incluyendo entornos de
integración continua sin hardware físico disponible. Todas las
\textit{fixtures} compartidas se definen de forma centralizada en
\texttt{tests/conftest.py}:

\begin{itemize}
    \item \texttt{mem\_db}: conexión SQLite en memoria (\texttt{:memory:})
    con el esquema completo inicializado. Empleada en los tests de base
    de datos puros, garantiza el aislamiento total entre ejecuciones.
    \item \texttt{test\_db}: ruta a un fichero \texttt{.db} temporal con
    todas las tablas creadas, útil para los tests que requieren una base
    de datos persistente durante el ciclo de vida del test.
    \item \texttt{mock\_serial}: parchea \texttt{serial.Serial}
    globalmente, simulando una conexión abierta con comportamientos
    configurables por cada test. Permite reproducir escenarios de
    desconexión, \textit{timeout} o recepción de tramas corruptas sin
    necesidad del hardware real.
    \item \texttt{mock\_reader}: instancia un \texttt{SerialReader} del
    paquete activo con la conexión serie ya vinculada al
    \texttt{mock\_serial}, eliminando la necesidad de reconfigurar
    el \textit{mock} en cada test individual.
    \item \texttt{mock\_model}: simula un modelo TensorFlow mediante
    \texttt{MagicMock}, devolviendo un vector de 5 probabilidades
    uniformes y la lista de las 5 clases reales en nuestro caso. Permite ejecutar el flujo
    completo de inferencia sin cargar TensorFlow en memoria.
    \item \texttt{window}: instancia completa de \texttt{SpectroControlUI}
    con la ruta de base de datos, el modelo y los puertos serie
    parcheados, registrada con \texttt{qtbot.addWidget()} para su
    limpieza automática al finalizar cada test.
\end{itemize}

Las dependencias que no se mockean son aquellas que ya incorporan
mecanismos de resiliencia propios: \texttt{inference/model.py} gestiona
internamente el \texttt{ImportError} de TensorFlow, y tanto
\texttt{widgets.py} como \texttt{pdf\_report.py} disponen de rutas de
\textit{fallback} ante la ausencia de \textbf{matplotlib}, por lo que
su comportamiento queda cubierto de forma indirecta por los tests
funcionales.

\subsection*{Pruebas de integración y hardware}

Dada la naturaleza del proyecto, las pruebas de integración se centran en
verificar la correcta comunicación entre el software y el hardware externo.
Se han realizado pruebas de estrés y validación sobre el puerto serie
(mediante \textbf{PySerial}) para garantizar que el sistema maneja
correctamente las desconexiones abruptas del Arduino Uno, los tiempos de
espera agotados (\textit{timeouts}) y la recepción de tramas de datos
incompletas o corruptas. Asimismo, se ha verificado la correcta interacción
entre la lógica de la aplicación y la base de datos SQLite, asegurando que
las operaciones de lectura y escritura sobre las tablas \texttt{muestras},
\texttt{predicciones}, \texttt{analisis} y \texttt{calibraciones} mantienen
la integridad referencial definida en el esquema.

\subsection*{Validación del modelo de IA y concurrencia}

Se han llevado a cabo pruebas específicas sobre el flujo de inferencia. Por
un lado, se ha validado que el modelo híbrido recibe los tensores con
la dimensionalidad y normalización correctas, emitiendo predicciones
consistentes con las métricas de precisión obtenidas durante su fase de
entrenamiento. La selección de esta arquitectura final estuvo precedida por una fase
de experimentación en la que se compararon múltiples algoritmos de
aprendizaje automático (Random Forest mediante \textbf{scikit-learn} y 
XGBoost mediante \textbf{xgboost}, entre otros), evaluando su rendimiento 
sobre el conjunto de datos espectrales recogido con el prototipo físico. 
Esta fase empírica justificó la adopción del modelo definitivo, cuyos resultados quedan recogidos 
en el capítulo de aspectos relevantes de la memoria principal.

Por otro lado, se ha comprobado exhaustivamente el comportamiento de la
interfaz gráfica durante el análisis concurrente, confirmando que la
delegación de tareas pesadas a hilos \textbf{QThread} evita bloqueos en
la GUI y permite una experiencia de usuario fluida y reactiva.

\subsection*{Análisis estático de calidad del código}

Con el objetivo de evaluar de forma objetiva la calidad del código
 desarrollado, se ha integrado la herramienta \textbf{SonarQube\cite{sonarqube2024docs}} en el
repositorio del proyecto mediante el fichero de configuración
\texttt{sonar\_project.properties}. La instancia utilizada es una instalación
autoalojada en \textbf{Docker}, desplegada en un servidor local, lo que permite mantener el análisis
dentro de la infraestructura local sin dependencias de servicios externos
en la nube. El análisis se lanza de forma manual desde el equipo de desarrollo
mediante \texttt{sonar-scanner}, autenticándose con el token de acceso personal
generado en dicha instancia.

Esta herramienta realiza un análisis estático del código, detectando
vulnerabilidades de seguridad, problemas de fiabilidad, deuda técnica de
mantenibilidad y duplicaciones. El análisis inicial del código fuente superó
el umbral de calidad (\textit{Quality Gate: Passed}), con los siguientes
resultados por categoría:

\begin{itemize}
    \item \textbf{Seguridad:} 0 incidencias abiertas — calificación \textbf{A}.
    \item \textbf{Fiabilidad:} 0 incidencias abiertas — calificación \textbf{A}.
    \item \textbf{Mantenibilidad:} 49 incidencias abiertas — calificación
    \textbf{A}. Estas incidencias corresponden principalmente a sugerencias
    de refactorización menores (\textit{code smells}) inherentes al carácter
    iterativo del desarrollo ágil, sin impacto funcional en el sistema.
    \item \textbf{Cobertura de pruebas:} 0{,}0\,\% sobre 8.800 líneas
    analizadas. SonarQube no recibía ningún informe de cobertura externo en
    esta configuración inicial, por lo que la métrica no refleja la batería
    de pruebas real del proyecto.
    \item \textbf{Duplicaciones:} 45{,}7\,\% sobre 8.800 líneas. Este índice
    elevado se explica por la coexistencia en el directorio \texttt{src/} de
    la codebase legada (\texttt{src/gui/}, \texttt{src/db/}) y el paquete
    activo refactorizado (\texttt{src/hives/}), que comparten estructuras
    equivalentes durante el periodo de transición arquitectónica.
    \item \textbf{Security Hotspots:} 0 — calificación \textbf{A}.
\end{itemize}

\imagentamano{img/Anexos/sonarqube_report}{Informe de calidad del código
generado por SonarQube en la configuración inicial (\texttt{sonar.sources=src}).
El proyecto supera el \textit{Quality Gate} con calificación A en seguridad,
fiabilidad y mantenibilidad; no obstante, la cobertura de pruebas registra
0{,}0\,\% por ausencia de informe de cobertura externo, y el índice de
duplicaciones alcanza el 45{,}7\,\% debido a la presencia simultánea de código
legado y código refactorizado en el repositorio. Fuente: Elaboración Propia.}{1.1}

\subsubsection*{Intentos de mejora de la integración con SonarQube}

Tras analizar los resultados obtenidos, se identificaron dos limitaciones en la
integración y se procedió a corregirlas experimentalmente.

\paragraph{Cobertura de pruebas.}
Se instrumentó la suite de pruebas con \texttt{pytest-cov} para generar un
informe en formato Cobertura XML:

\begin{verbatim}
pytest --cov=src/hives --cov-report=xml:coverage.xml tests/
\end{verbatim}

La ejecución confirmó una \textbf{cobertura real del 71{,}47\,\%} sobre 1.556
líneas del paquete activo (1.112 líneas cubiertas), con resultados destacados
en los módulos críticos: \texttt{reports/pdf\_report.py} (97{,}10\,\%),
\texttt{core/database.py} (93{,}75\,\%) e \texttt{inference/model.py}
(86{,}21\,\%). El fichero \texttt{coverage.xml} generado queda incluido en el
repositorio como evidencia.

Sin embargo, a pesar de añadir la propiedad \\
\texttt{sonar.python.coverage.reportPaths=coverage.xml} \\ al fichero de
configuración, la métrica de cobertura continuó mostrando 0{,}0\,\% en la
interfaz de SonarQube. La causa más probable es una incompatibilidad entre la
versión autoalojada de SonarQube y el formato del informe generado por la
versión de \texttt{coverage.py} utilizada.

\paragraph{Duplicaciones y restricción del alcance del análisis.}
Para reducir el índice de duplicaciones, se modificó la propiedad
\texttt{sonar.sources=src/hives} con el objetivo de restringir el análisis
únicamente al paquete activo. No obstante, la ejecución produjo un resultado
inesperado: el número de líneas analizadas se multiplicó desde las 8.800
iniciales hasta aproximadamente 212.000, y la calificación de
\textbf{Fiabilidad descendió a C}, con más de 12.000 incidencias reportadas.

Este comportamiento indica que la ruta especificada incluyó inadvertidamente
directorios de dependencias instaladas —presumiblemente el directorio
\texttt{\_internal} o el entorno virtual de Python presente en la misma
jerarquía—, los cuales fueron indexados por SonarQube junto con el código
propio. El elevado número de incidencias de fiabilidad es, por tanto,
consecuencia del análisis de bibliotecas de terceros, y no refleja la calidad
del código fuente desarrollado en el proyecto.

Ambas limitaciones quedan identificadas como líneas de mejora futura: la
correcta configuración de \texttt{sonar.exclusions} para excluir dependencias
y entornos virtuales, y la resolución de la incompatibilidad en el formato del
informe de cobertura.

\imagentamano{img/Anexos/sonarQubev2.png}{Informe de SonarQube tras la
modificación de \texttt{sonar.sources=src/hives}. El incremento no esperado de
líneas analizadas (hasta 212.000) y el descenso de Fiabilidad a C evidencian
que el analizador indexó dependencias externas incluidas en la ruta configurada,
y no únicamente el código fuente del proyecto. Fuente: Elaboración Propia.}{1}

\subsubsection*{Compilación y empaquetado del ejecutable}

La distribución de HIVES se realiza mediante \textbf{PyInstaller\cite{pyinstaller2024docs}}, herramienta
que genera un ejecutable autocontenido que encapsula el intérprete de Python,
todas las dependencias y el modelo de Inteligencia Artificial, eximiendo al
usuario final de cualquier instalación técnica previa. El proceso de empaquetado
queda íntegramente definido en el fichero \texttt{HIVES.spec}, ubicado en la
raíz del repositorio, que actúa como especificación reproducible del build.
Para generar el ejecutable basta con ejecutar desde la raíz del repositorio:

\begin{verbatim}
pyinstaller HIVES.spec
\end{verbatim}

El resultado se deposita en \texttt{dist/HIVES/}, siguiendo el formato
\textit{one-dir}: un ejecutable principal (\texttt{HIVES.exe} en Windows,
\texttt{HIVES} en Linux/macOS) acompañado de un directorio \texttt{\_internal/}
que contiene los binarios, las DLLs y el modelo \texttt{mejor\_modelo.keras}.
Este formato fue seleccionado frente al \textit{one-file} porque elimina la fase
de auto-extracción al arranque, reduciendo el tiempo de inicio de \mbox{4--8\,s}
a \mbox{2--4\,s}.

\paragraph{Optimizaciones aplicadas al \texttt{spec}.}
El fichero \texttt{HIVES.spec} incorpora un conjunto de optimizaciones
progresivas orientadas a reducir el tiempo de build y el tiempo de arranque
del ejecutable:

\begin{itemize}
    \item \textbf{Exclusión de paquetes no utilizados.} Se declaran en el
    parámetro \texttt{excludes} los paquetes que PyInstaller rastrearía
    innecesariamente: \texttt{pandas}, \texttt{scipy}, \texttt{sklearn},
    \texttt{cv2}, \texttt{setuptools}, \texttt{pdb}, \texttt{lib2to3},
    \texttt{curses}, \texttt{idlelib} y \texttt{turtledemo}, entre otros.
    La aplicación únicamente requiere PyQt6, TensorFlow, NumPy, matplotlib,
    FPDF2 y PySerial. Nótese que \texttt{PIL} y \texttt{unittest} no se
    excluyen porque matplotlib 3.10+ los referencia internamente. Esta
    medida reduce el tiempo de la fase de análisis de PyInstaller en
    aproximadamente un 20--40\,\%.

    \item \textbf{Optimización de bytecode.} El parámetro \texttt{optimize=2}
    instruye a Python para que compile los módulos embebidos eliminando
    \textit{docstrings} y aplicando optimizaciones de \textit{peephole}.
    Esto reduce ligeramente el tamaño de los \texttt{.pyc} y acelera la
    carga de módulos en torno a un 5--10\,\%, sin ningún impacto en la
    funcionalidad.

    \item \textbf{Desactivación de UPX.} La compresión UPX
    (\texttt{upx=False}) ha sido deshabilitada en Windows y Linux. Si bien
    UPX reduce el tamaño del bundle, añade tiempo de compresión en el build
    y tiempo de descompresión en cada arranque, lo que penaliza la experiencia
    de uso en el contexto de este proyecto.

    \item \textbf{Runtime hook de TensorFlow.} Se ha creado el fichero de \\
    \texttt{src/runtime\_hook.py}, referenciado en el parámetro
    \texttt{runtime\_hooks} del spec, que se ejecuta antes de cualquier
    import de la aplicación. Establece tres variables de entorno críticas:
    \texttt{TF\_CPP\_MIN\_LOG\_LEVEL=3} suprime los mensajes de diagnóstico
    de TensorFlow en consola;\\ \texttt{CUDA\_VISIBLE\_DEVICES=-1} impide la
    inicialización del subsistema CUDA/GPU, ahorrando entre 2 y 3 segundos
    de arranque en equipos sin tarjeta gráfica dedicada (perfil habitual del
    usuario objetivo) y evitando crashes por drivers ausentes;
    \texttt{OMP\_NUM\_THREADS=4} limita el paralelismo de OpenMP durante la
    carga de TensorFlow, evitando la proliferación de hilos competidores
    sobre la CPU.
\end{itemize}

\paragraph{Resolución del backend gráfico de matplotlib.}
Durante las pruebas del ejecutable empaquetado se detectó que el canvas
espectral mostraba el mensaje \textit{«matplotlib no disponible»}. La causa
raíz fue que PyInstaller no incluía el backend interactivo \texttt{QtAgg},
ya que su importación se produce dentro de un bloque \texttt{try/except} que
el analizador estático de PyInstaller no rastrea. La solución requirió tres
intervenciones coordinadas:

\begin{enumerate}
    \item Forzar la selección del backend antes de cualquier import,
    añadiendo\\ \texttt{matplotlib.use('QtAgg')} al inicio del runtime hook.
    \item Declarar explícitamente en \texttt{hiddenimports} los módulos \\
    \texttt{matplotlib.backends.backend\_qtagg}, \texttt{backend\_qt} y
    \texttt{backend\_agg}, así como sus dependencias transitivas
    (\texttt{kiwisolver}, \texttt{pyparsing},\\ \texttt{contourpy},
    \texttt{packaging}, \texttt{fonttools}, \texttt{dateutil}), completando
    el listado con \texttt{collect\_submodules('matplotlib.backends')}.
    \item Incluir los datos estáticos de matplotlib (fuentes y
    \texttt{matplotlibrc}) mediante \texttt{collect\_data\_files('matplotlib')}
    y un hook personalizado en \texttt{src/hooks/hook-matplotlib.py}.
\end{enumerate}

El conjunto de optimizaciones aplicadas reduce el tiempo de build de los
2--5 minutos originales a aproximadamente 30--60 segundos, y el tiempo de
arranque de 4--8 segundos a 2--4 segundos, manteniendo el 100\,\% de la
funcionalidad del sistema.
